\documentclass{beamer}

\mode<presentation>
\title{SoC1\&2: Tool Chain}
\author[AS]{Andreas Siggelkow \newline \texttt{andreas.siggelkow@rwu.de}}
\date[2026]{\today}
\subtitle{RWU}
\institute{Elektrotechnik und Informatik}

%%%%%%%%%%%%%%%%%%%%%%%%%%%%%%%%
% Logo
%%%%%%%%%%%%%%%%%%%%%%%%%%%%%%%%
\pgfdeclareimage[height=1cm]{hslogo}{./hs}
\logo{\pgfuseimage{hslogo}}
%%%%%%%%%%%%%%%%%%%%%%%%%%%%%%%%
% Macros
%%%%%%%%%%%%%%%%%%%%%%%%%%%%%%%%
\usepackage[T1]{fontenc}
\usepackage{lmodern}
\usepackage{textcomp}
\definecolor{links}{HTML}{2A1B81}
\hypersetup{colorlinks,linkcolor=,urlcolor=links}
\usepackage{graphicx}
\usepackage{ragged2e, tabularx}
\newcolumntype{Y}{>{\RaggedRight\arraybackslash}X}
\usetheme{Singapore}
\usecolortheme{seahorse}
\usefonttheme[onlylarge]{structuresmallcapsserif}
\usefonttheme[onlysmall]{structurebold}
\setbeamertemplate{footline}[frame number]
\usepackage{listings}
\newcommand{\astitle}{nix}
\newcommand{\assubtitle}{nix}
\newcommand{\asfolie}[3]{\begin{frame}\frametitle{#1}\begin{block}{#2}{#3}\end{block}\end{frame}}

% Commands and tool output are set with listings: dashes, $, _, \ and quotes
% appear exactly as they have to be typed.
\definecolor{lstbg}{HTML}{F2F2F7}
\definecolor{lstcomment}{HTML}{5A5A8A}
\lstset{
  basicstyle=\ttfamily\scriptsize,
  columns=fullflexible,
  keepspaces=true,
  upquote=true,
  breaklines=true,
  breakatwhitespace=false,
  backgroundcolor=\color{lstbg},
  frame=none,
  xleftmargin=2pt,
  xrightmargin=2pt,
  aboveskip=4pt,
  belowskip=4pt,
  showstringspaces=false
}
\lstdefinestyle{sh}{
  morecomment=[l]{\#},
  commentstyle=\color{lstcomment}\itshape
}
\lstdefinestyle{out}{basicstyle=\ttfamily\tiny}
\newcommand{\cmd}[1]{\texttt{#1}}

\begin{document}
%%%%%%%%%%%%%%%%%%%%%%%%%%%%%%%%
% First pages
%%%%%%%%%%%%%%%%%%%%%%%%%%%%%%%%
\setcounter{tocdepth}{1}

\begin{frame}
  \titlepage
\end{frame}

\begin{frame}
  \frametitle{Outline}
  \tableofcontents
\end{frame}

%%%%%%%%%%%%%%%%%%%%%%%%%%%%%%%%
% Tools
%%%%%%%%%%%%%%%%%%%%%%%%%%%%%%%%
\section{Tools}
\renewcommand{\astitle}{Tools}
\begin{frame}
\frametitle{\astitle}
\textcolor{red}{Tools}
\end{frame}

\renewcommand{\assubtitle}{Installation: RISC-V GNU Toolchain}
\begin{frame}[fragile]
\frametitle{\astitle}
\begin{block}{\assubtitle}
  \begin{itemize}
    \item Assembler, linker, objcopy
    \item Follow the current instructions on\\
      \url{https://github.com/riscv/riscv-gnu-toolchain}
    \item Recommended install target: \cmd{/opt/riscv}
    \item \cmd{make linux} builds the \cmd{riscv64-unknown-linux-gnu-*} tools,
      the prefix our CMake files use
  \end{itemize}
\begin{lstlisting}[style=sh]
git clone https://github.com/riscv/riscv-gnu-toolchain
cd riscv-gnu-toolchain
./configure --prefix=/opt/riscv
make linux
\end{lstlisting}
\end{block}
\end{frame}

\renewcommand{\assubtitle}{Installation: Caution with Vivado}
\begin{frame}
\frametitle{\astitle}
\begin{block}{\assubtitle}
  \begin{itemize}
    \item \textbf{Caution:} AMD Vivado ships its own RISC-V toolchain, which does not
      suit our needs. It may still be in your PATH (after sourcing \cmd{settings64.sh}).
    \item Our CMake files therefore call the tools with an absolute path:\\
      \cmd{\footnotesize /opt/riscv/bin/riscv64-unknown-linux-gnu-\{as,ld,objcopy\}}
  \end{itemize}
\end{block}
\end{frame}

\renewcommand{\assubtitle}{Installation: Simulator and Build System}
\begin{frame}
\frametitle{\astitle}
\begin{block}{\assubtitle}
  \begin{itemize}
    \item Simulator (one of them is enough)
    \begin{itemize}
      \item AMD Vivado simulator (xsim): \cmd{xvlog} / \cmd{xelab} / \cmd{xsim},
        waveforms in the Vivado GUI
      \item Verilator $\geq$ 5.0 (needed for \cmd{-{}-timing})
        plus GTKWave to view the FST waveform
    \end{itemize}
    \item Build system
    \begin{itemize}
      \item CMake $\geq$ 3.20 (plus make or ninja as the backend)
    \end{itemize}
  \end{itemize}
\end{block}
\end{frame}

%%%%%%%%%%%%%%%%%%%%%%%%%%%%%%%%
% Project Setup
%%%%%%%%%%%%%%%%%%%%%%%%%%%%%%%%
\section{Project Setup}
\renewcommand{\astitle}{Project Setup}
\begin{frame}
\frametitle{\astitle}
\textcolor{red}{Project Setup}
\end{frame}

\renewcommand{\assubtitle}{Project Layout (\texttt{RWU\_RV/RV\_NoPipeline})}
\begin{frame}
\frametitle{\astitle}
\begin{block}{\assubtitle}
  \small
  \begin{tabularx}{\linewidth}{@{}lY@{}}
    \cmd{rv64Sim/}        & assembler test programs (\cmd{*.asm}) + \cmd{CMakeLists.txt} (ASM $\rightarrow$ MEM)\\
    \cmd{src/}            & RTL of the CPU and the top level\\
    \cmd{../../ip\_*/}    & shared IPs (GPIO, UART, JTAG, CGU, DMem, ICU, trap CSR, \dots)\\
    \cmd{tb/}             & testbenches (\cmd{tb\_rv64i\_instr.sv} = shared instruction testbench)\\
    \cmd{sim/}            & \cmd{CMakeLists.txt} (compile + simulate), \cmd{sim\_main.cpp}, \cmd{*.tcl}, \cmd{*.wcfg}\\
    \cmd{CMakeLists.txt}  & top level: list of tests, testbench selection, regression\\
  \end{tabularx}
\end{block}
\end{frame}

\renewcommand{\assubtitle}{Three CMake Levels}
\begin{frame}[fragile]
\frametitle{\astitle}
\begin{block}{\assubtitle}
\begin{lstlisting}[basicstyle=\ttfamily\scriptsize]
CMakeLists.txt --> rv64Sim/CMakeLists.txt (target "mem")
               --> sim/CMakeLists.txt     (target "simulate"/"waves")
\end{lstlisting}
  \begin{itemize}
    \item The top level only orchestrates: one target per test
    \item \cmd{rv64Sim}: assembler program $\rightarrow$ memory image \cmd{rv64i.mem}
    \item \cmd{sim}: compile the RTL and run the simulation with that image
  \end{itemize}
\end{block}
\end{frame}

%%%%%%%%%%%%%%%%%%%%%%%%%%%%%%%%
% Single test case: simulation with Verilator and waveforms
%%%%%%%%%%%%%%%%%%%%%%%%%%%%%%%%
\section{Single Test Case}
\renewcommand{\astitle}{Single Test Case}
\begin{frame}
\frametitle{\astitle}
\textcolor{red}{Single test case: simulation with Verilator and waveforms}
\end{frame}

\renewcommand{\assubtitle}{Preparation}
\begin{frame}[fragile]
\frametitle{\astitle}
\begin{block}{\assubtitle}
\begin{lstlisting}[style=sh,basicstyle=\ttfamily\small]
cd RWU_RV/RV_NoPipeline
rm -rf build
cmake -S . -B build -DSIMULATOR=verilator -DWAVES=ON
cmake --build build --target sim_instr38ld
\end{lstlisting}
  \begin{itemize}
    \item \cmd{rm -rf build}: start from a clean build directory
    \item Two steps: \emph{configure} (\cmd{cmake -S \dots{} -B \dots}),
      then \emph{build} (\cmd{cmake -{}-build \dots})
  \end{itemize}
\end{block}
\end{frame}

\renewcommand{\assubtitle}{What the Options Mean}
\begin{frame}
\frametitle{\astitle}
\begin{block}{\assubtitle}
  \footnotesize
  \begin{tabularx}{\linewidth}{@{}lY@{}}
    \cmd{-S .}                        & source directory (the dot means the current directory)\\[2pt]
    \cmd{-B build}                    & build directory; it is created if needed and holds all generated files\\[2pt]
    \cmd{-DSIMULATOR=verilator}       & \cmd{-D} defines a CMake variable: simulator backend
                                        (\cmd{xsim} = default $|$ \cmd{verilator})\\[2pt]
    \cmd{-DWAVES=ON}                  & open the waveform viewer afterwards
                                        (single tests only, never together with \cmd{regression})\\[2pt]
    \cmd{-{}-build build}             & runs the build in the directory \cmd{build}\\[2pt]
    \cmd{-{}-target sim\_instr38ld}   & picks one test case: \cmd{sim\_<name of the .asm file>}\\
  \end{tabularx}
\end{block}
\end{frame}

\renewcommand{\assubtitle}{Under the Hood: What CMake Does}
\begin{frame}[fragile]
\frametitle{\astitle}
\begin{block}{\assubtitle}
  \begin{itemize}
    \item The top-level \cmd{CMakeLists.txt} defines one target \cmd{sim\_<test>} per test,
      each with its own build directory (here \cmd{build/instr38ld/}):
\begin{lstlisting}
set(TEST_BUILD_DIR ${CMAKE_BINARY_DIR}/${TEST})
\end{lstlisting}
    \item Testbench selection: \cmd{instr38ld} is listed in \cmd{INSTR\_TB\_CONFIG}
\begin{lstlisting}
"instr38ld"        "ld"          "38"
\end{lstlisting}
      $\rightarrow$ \cmd{tb/tb\_rv64i\_instr.sv} with \cmd{TEST\_NAME=ld},
      \cmd{EXPECTED=38} (GPIO value that signals a passing test)
    \item Tests that are not in that list use their own testbench
      \cmd{tb/tb\_rv64i\_<test>.sv}
  \end{itemize}
\end{block}
\end{frame}

\begin{frame}
\frametitle{\astitle}
\begin{block}{\assubtitle}
  \begin{itemize}
    \item Step 1: \cmd{build/instr38ld/asm/} is configured with
      \cmd{rv64Sim/CMakeLists.txt};\\
      the target \cmd{mem} produces \cmd{rv64i.mem}
    \item Step 2: \cmd{build/instr38ld/sim/} is configured with
      \cmd{sim/CMakeLists.txt};\\
      the target \cmd{waves} (\cmd{WAVES=ON}) or \cmd{simulate} (\cmd{WAVES=OFF}) runs
    \item Step 2 depends on \cmd{rv64i.mem}, so CMake always runs step 1 first
  \end{itemize}
\end{block}
\end{frame}

%%% ---------------- by hand ----------------
\renewcommand{\assubtitle}{The Same Steps by Hand (no CMake)}
\begin{frame}[fragile]
\frametitle{\astitle}
\begin{block}{\assubtitle}
  Shell variables used on the next slides (start in \cmd{RWU\_RV/RV\_NoPipeline}):
\begin{lstlisting}[style=sh,basicstyle=\ttfamily\small]
P=$PWD          # project root
IP=$P/../..     # IP root (ip_* directories)
RV=/opt/riscv/bin/riscv64-unknown-linux-gnu
mkdir -p build/instr38ld/asm build/instr38ld/sim
\end{lstlisting}
  \begin{enumerate}[a)]
    \item Assembler program $\rightarrow$ memory image
    \item Verilate and compile the C++ model
    \item Simulate
    \item Look at the waveforms
  \end{enumerate}
\end{block}
\end{frame}

\renewcommand{\assubtitle}{a) Assembler Program $\rightarrow$ Memory Image (1/2)}
\begin{frame}[fragile]
\frametitle{\astitle}
\begin{block}{\assubtitle}
\begin{lstlisting}[style=sh]
cd build/instr38ld/asm
$RV-as -march=rv64i_zicsr -mlittle-endian -g -o rv64i.o \
       $P/rv64Sim/instr38ld.asm -al=rv64i.lst
$RV-ld -b elf64-littleriscv -o rv64i.elf rv64i.o
\end{lstlisting}
  \begin{itemize}
    \item \cmd{as}: assemble; \cmd{-al=rv64i.lst} also writes a listing.
      CMake does not produce it, but it is handy for debugging.
    \item \cmd{ld}: link; an optional linker script can be passed with
      \cmd{-T <file>}, in CMake with \cmd{-DLINKER\_SCRIPT=<file>}
  \end{itemize}
\end{block}
\end{frame}

\renewcommand{\assubtitle}{a) Assembler Program $\rightarrow$ Memory Image (2/2)}
\begin{frame}[fragile]
\frametitle{\astitle}
\begin{block}{\assubtitle}
\begin{lstlisting}[style=sh]
$RV-objcopy -O verilog --verilog-data-width 4 rv64i.elf rv64i.mem
sed -i 's/ /\n/g' rv64i.mem   # one word per line
sed -i 's/\r//g'  rv64i.mem   # remove ^M (CR), if present
sed -i '1d'       rv64i.mem   # delete the address header
\end{lstlisting}
  \small
  \begin{itemize}
    \item \cmd{objcopy}: ELF $\rightarrow$ Verilog hex format,
      one 32-bit word = one instruction
    \item The linker places \cmd{.text} at \cmd{0x100B0}; the header holds the
      word address: \cmd{@0000402C} ($\mathtt{0x100B0} / 4 = \mathtt{0x402C}$)
    \item Without \cmd{sed '1d'}, \cmd{\$readmemh} would load the program there
      instead of at address 0, where the CPU starts. \alert{Required!}
  \end{itemize}
\end{block}
\end{frame}

\renewcommand{\assubtitle}{b) Verilate and Compile the C++ Model}
\begin{frame}[fragile]
\frametitle{\astitle}
\begin{block}{\assubtitle}
\begin{lstlisting}[style=sh,basicstyle=\ttfamily\tiny,breaklines=false]
cd ../sim
verilator --cc --exe $P/sim/sim_main.cpp --build --sv --timing \
  -Wall -Wno-fatal -Wno-DECLFILENAME -Wno-PINCONNECTEMPTY -Wno-IMPORTSTAR \
  --Mdir obj_dir -DLOGIC_SV --trace-fst --top-module tb_rv64i \
  $P/sim/tb_waivers.vlt \
  $P/src/as_pack.sv  $P/src/asRegFile.sv \
  $IP/ip_dmem/src/{asDMemBack,asDMemFront,asDMemCore,asDMem,asDMemTop}.sv \
  $P/src/asIMem.sv  $P/src/asAlu.sv  $P/src/asAluBr.sv  $P/src/asInstrDecode.sv \
  $IP/ip_jtag/src/{tap_fsm,ir_cell,dr_cell,scan_cell,by_cell,ir_reg,dr_reg,ir_decode,jtag}.sv \
  $P/src/asMBpi.sv  $P/src/asSBpi.sv \
  $IP/ip_gpio/src/asGpio.sv  $IP/ip_irqper/src/as_irq_peripheral.sv \
  $IP/ip_gpio/src/asGpioTop.sv \
  $IP/ip_trapcsr/src/cc_trap_csr.sv  $P/src/asCPUx.sv  $P/src/asDecode.sv \
  $IP/ip_cgu/src/{asCguCore,asCguTop}.sv \
  $IP/ip_fifo/src/as_fifo.sv  $IP/ip_uart_v2/src/as_uart_top.sv \
  $IP/ip_icu/src/{cc_icu,cc_icu_top}.sv \
  $P/src/asTopMem.sv \
  $P/tb/tb_rv64i_instr.sv
\end{lstlisting}
\end{block}
\end{frame}

\renewcommand{\assubtitle}{b) Verilator Options}
\begin{frame}
\frametitle{\astitle}
\begin{block}{\assubtitle}
  \small
  \begin{tabularx}{\linewidth}{@{}lY@{}}
    \cmd{-{}-cc -{}-exe \dots{} -{}-build} & generate C++ from the RTL, compile it together with
                                         \cmd{sim\_main.cpp} into \cmd{obj\_dir/Vtb\_rv64i}\\[2pt]
    \cmd{-{}-timing}     & supports the \cmd{\#}delays in the testbench\\[2pt]
    \cmd{-{}-trace-fst}  & \cmd{sim\_main.cpp} writes the waveform to \cmd{dump.fst}\\[2pt]
    \cmd{-DLOGIC\_SV}    & same define as for xsim (\cmd{xvlog -d LOGIC\_SV})\\
  \end{tabularx}
  \vspace{4pt}
  \begin{itemize}
    \item The order of the files matters: \cmd{as\_pack.sv} comes first,
      the testbench last
    \item The source list is in \cmd{sim/CMakeLists.txt} (\cmd{SOURCES\_SV})
  \end{itemize}
\end{block}
\end{frame}

\renewcommand{\assubtitle}{c) Simulate}
\begin{frame}[fragile]
\frametitle{\astitle}
\begin{block}{\assubtitle}
\begin{lstlisting}[style=sh]
cp ../asm/rv64i.mem riscvtest.mem
./obj_dir/Vtb_rv64i +TEST_NAME=ld +EXPECTED=38
\end{lstlisting}
  \begin{itemize}
    \item The file name is fixed: the testbench and \cmd{asIMem.sv} read
      \cmd{riscvtest.mem} with \cmd{\$readmemh} from the current working directory
    \item \cmd{+TEST\_NAME}, \cmd{+EXPECTED}: plusargs for \cmd{tb\_rv64i\_instr.sv}
    \item CMake does the same with:
  \end{itemize}
\begin{lstlisting}[style=sh]
cmake -E copy .../build/instr38ld/asm/rv64i.mem riscvtest.mem
cmake -P .../build/instr38ld/sim/run_sim.cmake
\end{lstlisting}
\end{block}
\end{frame}

\renewcommand{\assubtitle}{c) Expected Output and Evaluation}
\begin{frame}[fragile]
\frametitle{\astitle}
\begin{block}{\assubtitle}
\begin{lstlisting}
CS #1: gpio=0x26  expected=0x26
Simulation ld succeeded
%Error: .../tb/tb_rv64i_instr.sv:74: Verilog $stop
Aborting...
\end{lstlisting}
  \cmd{error.txt}: \cmd{11/10/26 - ld: Test ok}
  \footnotesize
  \begin{itemize}
    \item The testbench checks the GPIO value at every CS pulse:
    \begin{itemize}
      \item GPIO == \cmd{EXPECTED} $\rightarrow$ ``Simulation ld succeeded'',
        ``Test ok'' in \cmd{error.txt}
      \item no match after 20 CS pulses $\rightarrow$ ``Test fail'' in \cmd{error.txt}
    \end{itemize}
    \item The test ends with \cmd{\$stop}, so Verilator exits with code 1.
      That is not an error; \cmd{execute\_process} in \cmd{run\_sim.cmake} ignores it.
    \item \cmd{error.txt} is opened with \cmd{"a"}: results accumulate over runs!
  \end{itemize}
\end{block}
\end{frame}

\renewcommand{\assubtitle}{d) Look at the Waveforms}
\begin{frame}[fragile]
\frametitle{\astitle}
\begin{block}{\assubtitle}
\begin{lstlisting}[style=sh,basicstyle=\ttfamily\small]
gtkwave dump.fst &
\end{lstlisting}
  \begin{itemize}
    \item With \cmd{-DWAVES=ON}, CMake starts GTKWave the same way after the simulation
  \end{itemize}
\end{block}
\end{frame}

%%% ---------------- CMake output ----------------
\renewcommand{\assubtitle}{What CMake Prints: Configure}
\begin{frame}[fragile]
\frametitle{\astitle}
\begin{block}{\assubtitle}
\begin{lstlisting}[style=sh]
cmake -S . -B build -DSIMULATOR=verilator -DWAVES=ON
\end{lstlisting}
\begin{lstlisting}[style=out]
-- AS-TOP: WAVES=ON: Waveform GUI aktiv - nicht mit 'regression' kombinieren!
-- AS-TOP: SIMULATOR = verilator
-- AS-TOP: SIM_TARGET = waves
-- Configuring done (0.0s)
-- Generating done (0.0s)
-- Build files have been written to: .../RWU_RV/RV_NoPipeline/build
\end{lstlisting}
  \begin{itemize}
    \item The \cmd{AS-*} lines are our own messages from the \cmd{CMakeLists.txt} files
  \end{itemize}
\end{block}
\end{frame}

\renewcommand{\assubtitle}{What CMake Prints: Step 1, Configure the Assembler}
\begin{frame}[fragile]
\frametitle{\astitle}
\begin{block}{\assubtitle}
\begin{lstlisting}[style=sh]
cmake --build build --target sim_instr38ld
\end{lstlisting}
\begin{lstlisting}[style=out]
AS-TOP: Assembling instr38ld
 : cmake -E make_directory .../RWU_RV/RV_NoPipeline/build/instr38ld/asm
 : cmake -DASM_FILE=.../RWU_RV/RV_NoPipeline/rv64Sim/instr38ld.asm
 -S .../RWU_RV/RV_NoPipeline/rv64Sim -B .../RWU_RV/RV_NoPipeline/build/instr38ld/asm
 : cmake --build .../RWU_RV/RV_NoPipeline/build/instr38ld/asm --target mem
-- The ASM compiler identification is GNU
-- Found assembler: /usr/bin/cc
-- Configuring done (0.0s)
-- Generating done (0.0s)
-- Build files have been written to: .../RWU_RV/RV_NoPipeline/build/instr38ld/asm
\end{lstlisting}
  \begin{itemize}
    \item \cmd{/usr/bin/cc} is only CMake's language detection. The actual calls
      use \cmd{/opt/riscv/bin/...} (next slides).
  \end{itemize}
\end{block}
\end{frame}

\renewcommand{\assubtitle}{What CMake Prints: Step 1, Assemble and Link}
\begin{frame}[fragile]
\frametitle{\astitle}
\begin{block}{\assubtitle}
\begin{lstlisting}[style=out]
AS-rv64Sim: Assembling .../RWU_RV/RV_NoPipeline/rv64Sim/instr38ld.asm
 : /opt/riscv/bin/riscv64-unknown-linux-gnu-as
 -march=rv64i_zicsr
 -mlittle-endian
 -g
 -o .../RWU_RV/RV_NoPipeline/build/instr38ld/asm/rv64i.o
 .../RWU_RV/RV_NoPipeline/rv64Sim/instr38ld.asm

AS-rv64Sim: Linking ELF
 : /opt/riscv/bin/riscv64-unknown-linux-gnu-ld
 -b elf64-littleriscv
 -o .../RWU_RV/RV_NoPipeline/build/instr38ld/asm/rv64i.elf
 .../RWU_RV/RV_NoPipeline/build/instr38ld/asm/rv64i.o
\end{lstlisting}
\end{block}
\end{frame}

\renewcommand{\assubtitle}{What CMake Prints: Step 1, Memory Image}
\begin{frame}[fragile]
\frametitle{\astitle}
\begin{block}{\assubtitle}
\begin{lstlisting}[style=out]
AS-rv64Sim: Generating Verilog MEM
 : /opt/riscv/bin/riscv64-unknown-linux-gnu-objcopy
 -O verilog
 --verilog-data-width 4
 .../RWU_RV/RV_NoPipeline/build/instr38ld/asm/rv64i.elf
 .../RWU_RV/RV_NoPipeline/build/instr38ld/asm/rv64i.mem
 : sed -i 's/ /\n/g' .../RWU_RV/RV_NoPipeline/build/instr38ld/asm/rv64i.mem
 : sed -i 's/\r//g' .../RWU_RV/RV_NoPipeline/build/instr38ld/asm/rv64i.mem
 : sed -i '1d' .../RWU_RV/RV_NoPipeline/build/instr38ld/asm/rv64i.mem
[100%] Built target mem
\end{lstlisting}
\end{block}
\end{frame}

\renewcommand{\assubtitle}{What CMake Prints: Step 2, Configure the Simulation}
\begin{frame}[fragile]
\frametitle{\astitle}
\begin{block}{\assubtitle}
\begin{lstlisting}[style=out]
AS-TOP: Simulating instr38ld [verilator]
 : /usr/bin/cmake
 -DPROJECT_ROOT=.../RWU_RV/RV_NoPipeline
 -DIP_ROOT=.../RWU_RV/RV_NoPipeline/../..
 -DRV64_MEM=.../RWU_RV/RV_NoPipeline/build/instr38ld/asm/rv64i.mem
 -DTB_FILE=.../RWU_RV/RV_NoPipeline/tb/tb_rv64i_instr.sv
 -DSIMULATOR=verilator
 -DTB_TEST_NAME=ld;-DTB_EXPECTED=38
 -S .../RWU_RV/RV_NoPipeline/sim
 -B .../RWU_RV/RV_NoPipeline/build/instr38ld/sim
 : /usr/bin/cmake --build .../RWU_RV/RV_NoPipeline/build/instr38ld/sim --target waves
-- Configuring done (0.0s)
-- Generating done (0.0s)
-- Build files have been written to: .../RWU_RV/RV_NoPipeline/build/instr38ld/sim
\end{lstlisting}
  \begin{itemize}
    \item The \cmd{;} in \cmd{-DTB\_TEST\_NAME=ld;-DTB\_EXPECTED=38} is only how
      CMake prints a list: these are two separate arguments
  \end{itemize}
\end{block}
\end{frame}

\renewcommand{\assubtitle}{What CMake Prints: Step 2, Verilate (1/2)}
\begin{frame}[fragile]
\frametitle{\astitle}
\begin{block}{\assubtitle}
\begin{lstlisting}[style=out]
AS-sim: Verilating and building tb_rv64i
 : verilator
 --cc
 --exe .../RWU_RV/RV_NoPipeline/sim/sim_main.cpp
 --build
 --sv
 --timing
 -Wall
 -Wno-fatal
 -Wno-DECLFILENAME
 -Wno-PINCONNECTEMPTY
 -Wno-IMPORTSTAR
 --Mdir .../RWU_RV/RV_NoPipeline/build/instr38ld/sim/obj_dir
 -DLOGIC_SV
 --trace-fst
 --top-module tb_rv64i
 .../RWU_RV/RV_NoPipeline/sim/tb_waivers.vlt
 (source list: next slide)
\end{lstlisting}
\end{block}
\end{frame}

\renewcommand{\assubtitle}{What CMake Prints: Step 2, Verilate (2/2)}
\begin{frame}[fragile]
\frametitle{\astitle}
\begin{block}{\assubtitle}
  \cmd{SOURCES\_SV}, printed as one \cmd{;}-separated list
  (\cmd{./} = project root, \cmd{.../} = IP root):
\begin{lstlisting}[style=out]
./src/as_pack.sv;./src/asRegFile.sv;
.../ip_dmem/src/asDMemBack.sv;.../ip_dmem/src/asDMemFront.sv;
.../ip_dmem/src/asDMemCore.sv;.../ip_dmem/src/asDMem.sv;.../ip_dmem/src/asDMemTop.sv;
./src/asIMem.sv;./src/asAlu.sv;./src/asAluBr.sv;./src/asInstrDecode.sv;
.../ip_jtag/src/tap_fsm.sv;.../ip_jtag/src/ir_cell.sv;.../ip_jtag/src/dr_cell.sv;
.../ip_jtag/src/scan_cell.sv;.../ip_jtag/src/by_cell.sv;.../ip_jtag/src/ir_reg.sv;
.../ip_jtag/src/dr_reg.sv;.../ip_jtag/src/ir_decode.sv;.../ip_jtag/src/jtag.sv;
./src/asMBpi.sv;./src/asSBpi.sv;
.../ip_gpio/src/asGpio.sv;.../ip_irqper/src/as_irq_peripheral.sv;
.../ip_gpio/src/asGpioTop.sv;
.../ip_trapcsr/src/cc_trap_csr.sv;./src/asCPUx.sv;./src/asDecode.sv;
.../ip_cgu/src/asCguCore.sv;.../ip_cgu/src/asCguTop.sv;
.../ip_fifo/src/as_fifo.sv;.../ip_uart_v2/src/as_uart_top.sv;
.../ip_icu/src/cc_icu.sv;.../ip_icu/src/cc_icu_top.sv;
./src/asTopMem.sv;./tb/tb_rv64i_instr.sv
\end{lstlisting}
\end{block}
\end{frame}

\renewcommand{\assubtitle}{What CMake Prints: Step 2, Simulate and Waves}
\begin{frame}[fragile]
\frametitle{\astitle}
\begin{block}{\assubtitle}
\begin{lstlisting}[style=out]
AS-sim: Running Verilator simulation + GTKWave
 : cmake -E copy .../RWU_RV/RV_NoPipeline/build/instr38ld/asm/rv64i.mem riscvtest.mem
 : cmake -P .../RWU_RV/RV_NoPipeline/build/instr38ld/sim/run_sim.cmake
 : gtkwave dump.fst &
CS #1: gpio=0x26  expected=0x26
Simulation ld succeeded
%Error: .../RWU_RV/RV_NoPipeline/tb/tb_rv64i_instr.sv:74: Verilog $stop
Aborting...
\end{lstlisting}
  \begin{itemize}
    \item Same result as by hand: compare with ``c) Expected Output''
  \end{itemize}
\end{block}
\end{frame}

%%%%%%%%%%%%%%%%%%%%%%%%%%%%%%%%
% Variant: single test case with xsim
%%%%%%%%%%%%%%%%%%%%%%%%%%%%%%%%
\section{Single Test Case with xsim}
\renewcommand{\astitle}{Single Test Case with xsim}
\begin{frame}
\frametitle{\astitle}
\textcolor{red}{Variant: the same single test with xsim (Vivado)}
\end{frame}

\renewcommand{\assubtitle}{Preparation}
\begin{frame}[fragile]
\frametitle{\astitle}
\begin{block}{\assubtitle}
\begin{lstlisting}[style=sh,basicstyle=\ttfamily\small]
source <Vivado install dir>/settings64.sh
cd RWU_RV/RV_NoPipeline
rm -rf build
cmake -S . -B build -DSIMULATOR=xsim -DWAVES=ON
cmake --build build --target sim_instr38ld
\end{lstlisting}
  \begin{itemize}
    \item \cmd{settings64.sh} puts \cmd{xvlog}, \cmd{xelab} and \cmd{xsim} into the PATH
    \item \cmd{xsim} is the default: \cmd{-DSIMULATOR=xsim} may be left out
    \item Step 1 (assembler $\rightarrow$ \cmd{rv64i.mem}) is identical to Verilator;
      only step 2 in \cmd{sim/CMakeLists.txt} differs
  \end{itemize}
\end{block}
\end{frame}

\renewcommand{\assubtitle}{Under the Hood: Three xsim Steps}
\begin{frame}
\frametitle{\astitle}
\begin{block}{\assubtitle}
  \footnotesize
  \begin{tabularx}{\linewidth}{@{}lY@{}}
    \cmd{xvlog} & \emph{compile}: parses all SystemVerilog sources into the library \cmd{xsim.dir}\\[2pt]
    \cmd{xelab} & \emph{elaborate}: builds the design hierarchy below \cmd{tb\_rv64i}
                  and creates the snapshot \cmd{tb\_rv64i\_snapshot}\\[2pt]
    \cmd{xsim}  & \emph{simulate}: runs the snapshot, with or without the Vivado GUI\\
  \end{tabularx}
  \vspace{4pt}
  \begin{itemize}
    \item CMake targets in \cmd{sim/CMakeLists.txt}:
      \cmd{compile} $\rightarrow$ \cmd{elaborate} $\rightarrow$
      \cmd{simulate} (\cmd{WAVES=OFF}) or \cmd{waves} (\cmd{WAVES=ON})
    \item Each step leaves a stamp file (\cmd{compile.stamp}, \cmd{elaborate.stamp});
      CMake reruns a step only if the stamp is missing or older than its input
  \end{itemize}
\end{block}
\end{frame}

\renewcommand{\assubtitle}{The Same Steps by Hand: Compile and Elaborate}
\begin{frame}[fragile]
\frametitle{\astitle}
\begin{block}{\assubtitle}
  In \cmd{build/instr38ld/sim}, with \cmd{P} and \cmd{IP} as before:
\begin{lstlisting}[style=sh]
xvlog --sv --incr --relax -d LOGIC_SV <sources>
cp ../asm/rv64i.mem riscvtest.mem
xelab -debug all -top tb_rv64i -snapshot tb_rv64i_snapshot
\end{lstlisting}
  \footnotesize
  \begin{tabularx}{\linewidth}{@{}lY@{}}
    \cmd{<sources>}       & the same file list (and order) as for Verilator, without \cmd{tb\_waivers.vlt}\\[2pt]
    \cmd{-{}-sv}          & treat the files as SystemVerilog\\[2pt]
    \cmd{-{}-incr}        & incremental: only recompile files that changed\\[2pt]
    \cmd{-{}-relax}       & relaxed language checks\\[2pt]
    \cmd{-d LOGIC\_SV}    & the same define as \cmd{-DLOGIC\_SV} for Verilator\\[2pt]
    \cmd{-debug all}      & keep all signals visible for the waveform viewer\\
  \end{tabularx}
\end{block}
\end{frame}

\renewcommand{\assubtitle}{The Same Steps by Hand: Simulate}
\begin{frame}[fragile]
\frametitle{\astitle}
\begin{block}{\assubtitle}
  With the Vivado GUI and waveforms (\cmd{WAVES=ON}):
\begin{lstlisting}[style=sh]
xsim --gui tb_rv64i_snapshot \
     --testplusarg "TEST_NAME=ld" --testplusarg "EXPECTED=38" \
     --tclbatch $P/sim/xsim_cfg2.tcl
\end{lstlisting}
  Without the GUI (\cmd{WAVES=OFF}):
\begin{lstlisting}[style=sh]
xsim tb_rv64i_snapshot \
     --testplusarg "TEST_NAME=ld" --testplusarg "EXPECTED=38" \
     --tclbatch $P/sim/xsim_cfg.tcl
\end{lstlisting}
  \begin{itemize}
    \item \cmd{-{}-testplusarg} passes the same plusargs as
      \cmd{+TEST\_NAME} / \cmd{+EXPECTED} for Verilator
  \end{itemize}
\end{block}
\end{frame}

\renewcommand{\assubtitle}{The Tcl Scripts}
\begin{frame}[fragile]
\frametitle{\astitle}
\begin{block}{\assubtitle}
  \begin{columns}[T]
  \column{0.48\textwidth}
  \cmd{sim/xsim\_cfg.tcl} (batch)
\begin{lstlisting}
log_wave -recursive * -v
run all
exit
\end{lstlisting}
  \column{0.48\textwidth}
  \cmd{sim/xsim\_cfg2.tcl} (GUI)
\begin{lstlisting}
log_wave -recursive *
add_wave {{*}}
run all
#exit
\end{lstlisting}
  \end{columns}
  \footnotesize
  \begin{itemize}
    \item \cmd{run all} runs until the testbench's \cmd{\$stop}
    \item Batch: \cmd{exit} ends xsim afterwards
    \item GUI: no \cmd{exit}, so the Vivado window stays open with all
      top-level signals in the waveform; add more signals from the scope tree
      (or load one of the \cmd{sim/*.wcfg} files)
    \item The result is evaluated as with Verilator: console output and
      \cmd{error.txt} in \cmd{build/instr38ld/sim}
  \end{itemize}
\end{block}
\end{frame}

%%%%%%%%%%%%%%%%%%%%%%%%%%%%%%%%
% Regression
%%%%%%%%%%%%%%%%%%%%%%%%%%%%%%%%
\section{Regression}
\renewcommand{\astitle}{Regression}
\begin{frame}
\frametitle{\astitle}
\textcolor{red}{Regression: all tests at once}
\end{frame}

\renewcommand{\assubtitle}{Run All Tests}
\begin{frame}[fragile]
\frametitle{\astitle}
\begin{block}{\assubtitle}
\begin{lstlisting}[style=sh,basicstyle=\ttfamily\small]
cd RWU_RV/RV_NoPipeline
rm -rf build
cmake -S . -B build -DWAVES=OFF        # xsim (default)
cmake --build build --target collect_errors -j8
\end{lstlisting}
  \begin{itemize}
    \item With Verilator: add \cmd{-DSIMULATOR=verilator} as for a single test
    \item \cmd{-DWAVES=OFF} is required: otherwise every test opens a waveform viewer
    \item \cmd{-j8}: up to 8 tests run in parallel. That is safe, because every test
      has its own directory \cmd{build/<test>/}
    \item Result: \cmd{regression.txt} in the project root
  \end{itemize}
\end{block}
\end{frame}

\renewcommand{\assubtitle}{What Happens}
\begin{frame}[fragile]
\frametitle{\astitle}
\begin{block}{\assubtitle}
\begin{lstlisting}[basicstyle=\ttfamily\scriptsize]
collect_errors
  `-- depends on: regression
        `-- depends on: sim_<test>   for every test in ASM_TESTS
              (each one exactly as in "Single Test Case")
  then: cmake -P collect_errors.cmake
        build/<test>/sim/error.txt  -->  regression.txt
\end{lstlisting}
  \begin{itemize}
    \item \cmd{ASM\_TESTS} in the top-level \cmd{CMakeLists.txt} lists all tests
      (currently 64)
    \item \cmd{regression} alone only runs the tests; \cmd{collect\_errors}
      also writes the report
    \item \cmd{regression.txt} is rewritten on every run;
      the \cmd{error.txt} files are not (see Pitfalls)
  \end{itemize}
\end{block}
\end{frame}

\renewcommand{\assubtitle}{Reading the Report}
\begin{frame}[fragile]
\frametitle{\astitle}
\begin{block}{\assubtitle}
  \begin{columns}[T]
  \column{0.5\textwidth}
\begin{lstlisting}[basicstyle=\ttfamily\tiny]
=== RV64I REGRESSION REPORT ===

--- instr99scantest ---
11/06/26 - scan test: Test ok

--- instr01loadbyte ---
11/06/26 - load_byte: Test ok

--- instr38ld ---
11/06/26 - ld: Test ok
...
\end{lstlisting}
  \column{0.48\textwidth}
  \footnotesize
  One section per test with the content of its \cmd{error.txt}:
  \begin{itemize}
    \item \cmd{Test ok}: passed
    \item \cmd{Test fail}: failed
    \item \cmd{MISSING error.txt}: the test did not
      report at all. Treat it as a failure!
  \end{itemize}
  \end{columns}
  \vspace{4pt}
  Quick check:
\begin{lstlisting}[style=sh]
grep -c "Test ok" regression.txt               # 64 expected
grep -B1 "fail\|MISSING" regression.txt           # which tests?
\end{lstlisting}
\end{block}
\end{frame}

\renewcommand{\assubtitle}{Adding a New Test}
\begin{frame}[fragile]
\frametitle{\astitle}
\begin{block}{\assubtitle}
  \begin{enumerate}
    \item Write the program \cmd{rv64Sim/<test>.asm}
    \item Add \cmd{<test>} to \cmd{ASM\_TESTS} in the top-level \cmd{CMakeLists.txt}
    \item Choose the testbench:
    \begin{itemize}
      \item shared instruction testbench: add a line to \cmd{INSTR\_TB\_CONFIG}
\begin{lstlisting}
"<test>"   "<name in error.txt>"   "<expected GPIO value>"
\end{lstlisting}
      \item otherwise: write your own \cmd{tb/tb\_rv64i\_<test>.sv}
    \end{itemize}
    \item Configure again, then test it alone first:
\begin{lstlisting}[style=sh]
cmake -S . -B build
cmake --build build --target sim_<test>
\end{lstlisting}
  \end{enumerate}
\end{block}
\end{frame}

%%%%%%%%%%%%%%%%%%%%%%%%%%%%%%%%
% Pitfalls
%%%%%%%%%%%%%%%%%%%%%%%%%%%%%%%%
\section{Pitfalls}
\renewcommand{\astitle}{Pitfalls}
\renewcommand{\assubtitle}{Toolchain and Cache Variables}
\begin{frame}[fragile]
\frametitle{\astitle}
\begin{block}{\assubtitle}
  \begin{itemize}
    \item \textbf{Vivado's RISC-V toolchain in the PATH}\\
      After \cmd{source settings64.sh}, \cmd{riscv64-unknown-*} may resolve to Vivado's tools.
      Always call \cmd{/opt/riscv/bin/...} with the absolute path.
    \item \textbf{\cmd{WAVES} and \cmd{SIMULATOR} are cache variables}\\
      They stay set in \cmd{build/CMakeCache.txt} until you change them explicitly
      or delete \cmd{build/}:
  \end{itemize}
\begin{lstlisting}[style=sh]
cmake -S . -B build -DWAVES=OFF    # before every regression!
\end{lstlisting}
  \begin{itemize}
    \item Never run a regression with \cmd{WAVES=ON}: one GUI per test
  \end{itemize}
\end{block}
\end{frame}

\renewcommand{\assubtitle}{RTL Changes Are Not Detected}
\begin{frame}[fragile]
\frametitle{\astitle}
\begin{block}{\assubtitle}
  \begin{itemize}
    \item The simulation step only depends on \cmd{rv64i.mem}, and \cmd{rv64i.mem}
      only on the \cmd{.asm} file
    \item After a change to the RTL (\cmd{.sv}), CMake rebuilds \emph{nothing}:
      the old model is simulated again!
    \item Applies to both simulators (\cmd{obj\_dir/Vtb\_rv64i} and \cmd{compile.stamp})
  \end{itemize}
  Before the next run:
\begin{lstlisting}[style=sh]
rm -rf build                     # everything, or at least
rm -rf build/instr38ld/sim       # the simulation part of one test
\end{lstlisting}
\end{block}
\end{frame}

\renewcommand{\assubtitle}{\texttt{error.txt} Accumulates}
\begin{frame}[fragile]
\frametitle{\astitle}
\begin{block}{\assubtitle}
  \begin{itemize}
    \item The testbench opens \cmd{error.txt} with \cmd{"a"} (append):
      every run adds a line, old results stay
\begin{lstlisting}
11/10/26 - ld: Test fail
11/12/26 - ld: Test ok
\end{lstlisting}
    \item \cmd{regression.txt} copies these files completely,
      so it also shows old results
    \item Delete old results before you judge them:
\begin{lstlisting}[style=sh]
rm -f build/*/sim/error.txt
\end{lstlisting}
    \item Also: Verilator's exit code 1 after \cmd{\$stop} is \emph{not} a failure;
      only \cmd{error.txt} counts
  \end{itemize}
\end{block}
\end{frame}

%%%%%%%%%%%%%%%%%%%%%%%%%%%%%%%%
% Q&A
%%%%%%%%%%%%%%%%%%%%%%%%%%%%%%%%
\section{Q\&A}
\renewcommand{\astitle}{Q\&A}
\begin{frame}
\frametitle{\astitle}
\vfill
\centering
\begin{Huge}\textcolor{red}{Questions?}\end{Huge}
\vfill
\end{frame}

\end{document}
